\documentclass[11pt,a4paper]{article}
\usepackage[margin=25mm]{geometry}
\usepackage{amsmath,amssymb,amsthm,hyperref}
\setlength{\emergencystretch}{2em}
\newtheorem{lemma}{Lemma}\newtheorem{theorem}[lemma]{Theorem}
\newcommand{\gen}[1]{\langle#1\rangle}
\newcommand{\Sym}{\operatorname{Sym}}
\title{A first split away from the anchor}
\author{Human Compression project: a component of the L019 recovery}
\date{}
\begin{document}\maketitle
We retain the actual critical configuration of the L019 recovery.
The full marking is a fractured natural symmetric string C-marking
on $X\cup Y$, where $|X|=m$ is even, $|Y|=6$, and
$X\cap Y=\{e\}$. Its labels are the concatenation of an even
degree-six rank-three $A_5$ factor $E$ on $Y$ and a natural
symmetric factor
\[
                 P=\gen{t_0,\ldots,t_{s-1}}=\Sym(X),
                 \qquad s=m/2+1.
\]
The common point $e$ is an anchor at the last mark of $E$ and
the first mark of $P$. Both factors are fractured. The full
degree and rank are $n=m+5$ and $r=m/2+4$, so $2r=n+3$.

An internal perfect split is called \emph{sufficient} if its
split mark is pure, or if one internal factor is trivial and
the two actual groups meeting in the crossing point both
contain an odd permutation. The acquired pure-split and
odd-attachment theorems give the required subdivision in
these cases. In particular, when the split mark is
$q=(c\ d)\beta$ and its supported wings are $F,R$, these
two groups are $F$ and $\gen{q,R}$, not merely $F,R$.

Assume throughout the strict smaller-degree sufficient-split
induction: every fractured natural symmetric string C-marking
of degree $d<n$ and rank $u$ with $2u>d+2$ has an internal
sufficient split. No critical factor at equality is assumed
solved by this hypothesis.

\begin{lemma}\label{right}
Let $j$ be the least colour of $P$ with exactly one eligible
edge. It is a perfect split. Write $A,B$ for its omission
orbits, with $e\in A$, and put $a=|A|$, $b=|B|$.
Then $2j\le a-1$, the right wing is the actual $\Sym(B)$,
and $j\ne1$.
\end{lemma}
\begin{proof}
The acquired anchored first-split theorem applies to $P$:
it is transitive and fractured, has the given anchor, and
$2s=m+2>m$. It supplies perfection and the stated density.
The right wing has rank $q=s-j-1=m/2-j$, whence
\[
                         2q=b+a-2j\ge b+1.             \tag{1}
\]
Use the precise nonsymmetric-wing conclusion of the acquired
actual-wing theorem. It bounds each actual nonsymmetric wing
by half its degree, independently of the stronger parent-rank
inequality used only in that theorem's final either-wing
conclusion. Its full L047 exception list and its $A_{11}$
attachment obstruction are retained inputs here. Equation (1)
therefore forces the right wing to be $\Sym(B)$.
If $j=1$, its single left involution would be transitive
on $A$, so $a=2$, contrary to $2j\le a-1$.
\end{proof}

\begin{lemma}\label{transport}
An internal sufficient split in an actual symmetric wing
with $2u>d+2$ transfers to an internal sufficient split of
the full marking. The same holds for a symmetric factor
supported on one of two sets meeting in one point, with
consecutive label intervals and the stated endpoint anchors.
\end{lemma}
\begin{proof}
For a wing use the acquired internal split-inheritance proof.
An internal index has distance at least two from the old
split mark. That mark commutes with the relevant transitive
side and with the new unique crossing, and hence fixes that
side pointwise. The opposite side becomes its old side
together with the other wing, and remains transitive.
The new crossing and its internal factors are unchanged.
Purity and a trivial factor therefore persist. Each new
attaching group contains its corresponding old attaching
group, so each old odd element persists as an actual
permutation. This also preserves sufficiency.

For factors meeting in one point, use the acquired one-point
transport lemma: the side containing the shared anchor is
enlarged by the other factor, and the other side is fixed.
The same containment and parity argument applies. All such
submarkings are fractured, since a transitive omission there
would become a transitive omission of the full marking.
Their restrictions to their actual supports are faithful
and inherit IP; no quotient is used.
\end{proof}

\begin{theorem}\label{away}
Under the standing induction hypothesis, if $j>0$ in
Lemma~\ref{right}, the full marking has an internal
sufficient split. Consequently, a remaining critical
configuration may be assumed to have $j=0$.
\end{theorem}
\begin{proof}
By Lemma~\ref{right}, $j\ge2$ and $a\ge2j+1\ge5$.
The split of $P$ is also a perfect split of the full
marking. Its left wing
\[
 F=\gen{E,t_0,\ldots,t_{j-1}}
\]
is supported and transitive on $Y\cup A$, of degree
$a+5$ and rank $j+3$. Indeed $E$ and the prefix are
transitive on sets meeting at $e$. The right wing
$R=\Sym(B)$ is unchanged. The split mark $t_j$ fixes
$Y$ pointwise, and has the form
\[
                  t_j=(c\ d)\alpha\beta,
   \qquad c\in A,\quad d\in B,
\]
with the two internal factors on their respective sides.
Set $\delta=a-2j\ge1$. Equation (1) becomes
\[
                             2q=b+\delta.              \tag{2}
\]
The right wing is nontrivial: otherwise $b=1,q=0$
would contradict (2).

Suppose first that both internal factors are nontrivial.
Apply the acquired faithful-augmentation lemma with $R$
as its symmetric wing, reversing the labels. The other
augmented factor has a faithful primitive C-action $K$
of degree $a+6$ and rank $j+4$. That lemma excludes both
natural alternating and symmetric groups and supplies
its nontrivial first-mark character. If $\delta\le2$,
\[
                      2(j+4)=(a+6)+2-\delta\ge a+6.
\]
The complete primitive theorem of L047 now applies at
or above half the degree. Its proper exceptions have
degrees six and ten, while $a+6\ge11$. Thus this case
is impossible. It follows that $\delta\ge3$.
Equation (2) gives $2q>b+2$, so induction in the smaller
actual symmetric wing $R$, followed by
Lemma~\ref{transport}, proves the assertion.

Suppose next that $\alpha=1$. Since $\beta\in R$,
the group $J=\gen{t_j,R}$ is the actual
$\Sym(B\cup\{c\})$. It is supported on that set,+meeting the support of $F$ only at $c$. If $F$ is odd,
both attaching groups are odd and this split is sufficient.
If $F$ is even and $\delta\ge2$, then
\[
              2\operatorname{rank}(J)=b+\delta+2
                            \ge |B\cup\{c\}|+3.
\]
Induction in $J$ and one-point transport apply, since
its degree $b+1<n$.

It remains in this case that $F$ is even and $\delta=1$.
Then its degree $f=a+5$ is even, $f\ge10$, and its
rank is $f/2$. It is fractured and has an endpoint
anchor at $c$. The acquired anchored-even theorem
excludes $f\ge12$. At $f=10$, the acquired small-even
theorem applies with the actual odd external group $J$.
Its only permissible even half-degree factor is the
degree-six $A_5$, so degree ten is excluded too.
These invocations retain the exact actual attachment
and supported-kernel hypotheses supplied by $F,J$.

Finally suppose $\beta=1$; the case where both factors
are trivial already gives a pure split. Reverse the
split. The two factors meeting at $d$ are now $R$ and
$H=\gen{F,t_j}$, supported on $B$ and
$Y\cup A\cup\{d\}$. The symmetric group $R$ is odd.
If $H$ is odd the split is sufficient. If $H$ is even,
the acquired even-attachment theorem says either that
$2q>b+2$, or that $H$ is the degree-six $A_5$.
The latter is impossible since $\deg(H)=a+6\ge11$.
In the former case induction in $R$ and wing transport
finish the proof.

Every induction above uses strictly fewer than $n$
points and the strict rank inequality. If a resulting
degree were at most four, that strict inequality is
already impossible: a fractured transitive family on
$d$ points has at most $d-1$ members by the acquired
one-edge orbit forest. Thus no unstated small base or
critical equality is being invoked.
\end{proof}

\paragraph{Attribution and remaining case.}
The rank classifications, actual-wing theorem, faithful
augmentation, even-attachment and finite even-factor
exceptions remain fully attributed in the acquired
components. The present deduction combines them with
the new first-cut density. It does not make an arbitrary
central quotient into a C-group. If $j=0$, the symmetric
factor begins with its unique crossing at its anchor,
but that mark may still have a nontrivial internal
factor. The remaining endpoint configuration and the
stable counting correspondence remain unproved here.
No whole-original, conjecture, Lean or measurement
credit is asserted.
\end{document}
